%------------------------------ Ejercicio 6.23 -------------------------------
\setcounter{neje}{22}
\begin{nejer}
Fije dos números reales $a < b$ y pruebe lo siguiente.
\begin{enumerate}[label = (\roman*)]
    \item Si $f \colon [a, b] \to \mathbb{R}$ es diferenciable, entonces $f' \colon [a, b] \to \mathbb{R}$ es Borel medible.
    \item Si $f \colon [a, b] \to \mathbb{R}$ es diferenciable y $\int_a^b |f'(t)| \dd{t} < \infty$, entonces $f$ es absolutamente continua.
\end{enumerate}
\emph{Sugerencia.} Fije una constante $\varepsilon > 0$. Por el teorema de Vitali-Carathéodory en el Ejercicio 6.22, existe una función semi-continua inferiormente $g \colon [a, b] \to \mathbb{R}$ tal que
$$
g > f', \quad \int_a^b g(t) \dd{t} < \int_a^b f'(t) \dd{t} + \varepsilon.
$$
Para $\eta > 0$ defina la función
$$
F_\eta \colon [a, b] \to \mathbb{R}
$$
mediante
$$
F_\eta(x) := \int_a^x g(t) \dd{t} - f(x) + f(a) + \eta(x - a)
$$
para $a \le x \le b$. Considere un punto $a \le x < b$. Puesto que $g(x) > f'(x)$ y $g$ es semi-continua inferiormente, encuentre un número $\delta_x > 0$ tal que
$$
g(t) > f'(x), \quad \frac{f(t) - f(x)}{t - x} < f'(x) + \eta \quad \text{para } x < t < x + \delta_x.
$$
Deduzca que
$$
F_\eta(t) > F_\eta(x) \quad \text{para } x < t < x + \delta_x.
$$
Puesto que $F_\eta(a) = 0$, existe un elemento maximal $x \en [a, b]$ tal que $F_\eta(x) = 0$. Si $x < b$, se sigue de la discusión anterior que $F_\eta(t) > 0$ para $x < t \le b$. En cualquier caso $F_\eta(b) \ge 0$ y por lo tanto
$$
f(b) - f(a) \le \int_a^b g(t) \dd{t} + \eta(b - a) < \int_a^b f'(t) \dd{t} + \varepsilon + \eta(b - a).
$$
Puesto que esto se cumple para todo $\eta > 0$ y todo $\varepsilon > 0$, se sigue que
$$
f(b) - f(a) \le \int_a^b f'(t) \dd{t}.
$$
Reemplace $f$ por $-f$ para obtener la relación $f(b) - f(a) = \int_a^b f'(t) \dd{t}$. Ahora deduzca que
$$
f(x) - f(a) = \int_a^x f'(t) \dd{t}
$$
para todo $x \en [a, b]$.
\end{nejer}
\begin{sol}
\begin{enumerate}[label = (\roman*)]
\item Sea $f \colon [a,b]\to\mathbb{R}$ diferenciable. Para cada $n\en\mathbb{N}$ defina
\[
\varphi_n(x)
\coloneq
n\bigl(f(x+\tfrac1n)-f(x)\bigr),
\qquad\forall x\en[a,b-\nicefrac{1}{n}]
\]
Podemos extender $\varphi_n$ al intervalo completo $[a,b]$ haciendo que sea nula en $(b-\nicefrac{1}{n},b]$, así será continua por trozos y por lo tanto Borel-medible.
Como $f$ es diferenciable,
\[
f'(x)
=
\lim_{n\to\infty}
n\bigl(f(x+\tfrac1n)-f(x)\bigr)
= \lim_{n\to\infty} \varphi_n(x)
\]
para todo $x\en[a,b)$. En consecuencia, $f'$ es el límite puntual de funciones Borel-medibles, y por lo tanto es Borel-medible.

\item Por el inciso anterior $f'$ es medible, ahora supongamos $\int_a^b |f'(t)|\,\dd t < \infty$, i.e., $f'\en\mathcal{L}^1([a,b])$.
Fijemos $\varepsilon>0$. Por el teorema de Vitali--Carathéodory (Ejercicio 6.22, tarea 9) existe una función semi-continua y acotada inferiormente $g\colon [a,b]\to\mathbb{R}$ tal que $f' < g$
\[
\int_a^b g(t)\,\dd t
<
\int_a^b f'(t)\,\dd t+\varepsilon.
\]
Para cada $\eta>0$ definimos $F_\eta \colon [a,b] \to \mathbb{R}$ tal que
\[
F_\eta(x)
:=
\int_a^x g(t)\,\dd t
-
f(x)
+
f(a)
+
\eta(x-a),
\]
y fijemos $x\en[a,b)$. Como $g(x)>f'(x)$ y $g$ es semi-continua inf., existe $\delta_1>0$ tal que
\[
t\en(x,x+\delta_1)
\quad\implies\quad
g(t)>f'(x)
\]
Además, gracias a que $f'(x)=\lim_{t\to x^+}\frac{f(t)-f(x)}{t-x}$, existe $\delta_2>0$ tal que
\[
t\en(x,x+\delta_2)
\quad\implies\quad
\frac{f(t)-f(x)}{t-x}
<
f'(x)+\eta
\]
Si tomamos $\delta_x:=\min\{\delta_1,\delta_2\}$ entonces para todo $t \en (x,x+\delta_x)$
\[
F_\eta(t)-F_\eta(x)
=
\int_x^t g(s)\,\dd s
-
\bigl(f(t)-f(x)\bigr)
+
\eta(t-x).
\]
Como $g(s)>f'(x)$ para $s\en(x,t) \subseteq (x,x+\delta_1)$,
\[
    \int_x^t g(s)\,\dd s
    >
    f'(x)(t-x),
\]
y a su vez, gracias a que $t \en (x,x+\delta_x) \subseteq (x,x+\delta_2)$:
\[
    f(t)-f(x)
    <
    (f'(x)+\eta)(t-x).
\]
Por lo tanto,
{\thickmuskip=10mu
\[
F_\eta(t)-F_\eta(x)
>
f'(x)(t-x)
-
(f'(x)+\eta)(t-x)
+
\eta(t-x)
=0
\]
}%
Es decir, hemos probado que $F_\eta(t)>F_\eta(x)$ para todo $t\en(x,x+\delta_x)$.

Ahora bien, por la implicación "(iii) $\Rightarrow$ (i)" del teorema 6.19 (TFC), sabemos que el mapeo $x \mapsto \int^x_ag(t)\dd{t}$ es absolutamente continuo, pero $f(x)$ también es continua al ser diferenciable. De esto se sigue que $F_\eta(x)$ es continua, por lo cual
\[
E \coloneq F_\eta^{-1}(\set0) \subseteq [a,b]
\]
es compacto. Como además $E\neq\varnothing$ gracias a que $F_\eta(a) =0$, existe $x_0\coloneq\max E \en [a,b]$. Ahora afirmamos que $F_\eta(b) \ge0$. Si $x_0 = b$ ya terminamos, puesto que $F_\eta(b) = 0$. Si $x_0<b$, la discusión anterior prueba que
\[
t\en(x_0,x_0+\delta_{x_0})
\quad\implies\quad
F_\eta(t)>F_\eta(x_0)=0
\]
Así, la ausencia de ceros en $(x_0,b]$ (maximalidad de $x_0$), fuerza a que $F_\eta(b) \ge 0$.%
\footnote{
De otro modo, $F_\eta(b) < 0$ y $F(x_0+\frac12\delta_{x_0}) >0$ implicaría, por continuidad de $F_\eta$, que existe $t \en (x_0+\frac12\delta_{x_0}, b)$ con $F_\eta(t) = 0$, lo que contradice la maximalidad de $x_0$.
}
Entonces
\[
    \int_a^b g(t)\,\dd t
    -
    f(b)
    +
    f(a)
    +
    \eta(b-a)
    \ge 0,
\]
es decir,
\[
    f(b)-f(a)
    \le
    \int_a^b g(t)\,\dd t
    +
    \eta(b-a).
\]
Recordando que $\int g(t) \dd{t} < \int f'(t)\dd{t} + \varepsilon$:
\[
    f(b)-f(a)
    <
    \int_a^b f'(t)\,\dd t
    +
    \varepsilon
    +
    \eta(b-a).
\]
Como esto se cumple para cualesquiera $\varepsilon>0$ y $\eta>0$, concluimos que
\[
    f(b)-f(a)
    \le
    \int_a^b f'(t)\,\dd t.
\]

Aplicando el mismo razonamiento a $-f$, obtenemos
\[
    -(f(b)-f(a))
    \le
    -\int_a^b f'(t)\,\dd t,
\]
y por lo tanto
\[
    f(b)-f(a)
    =
    \int_a^b f'(t)\,\dd t.
\]
Finalmente aplicando lo recién probado a $f|_{[a,x]}$ para $x \en [a,b]$, se sigue
\[
    f(x)-f(a)
    =
    \int_a^x f'(t)\,\dd t
\]
Por último, la implicación "(iii) $\Rightarrow$ (i)" del teorema 6.19 garantiza que $f$ es absolutamente continua.
\end{enumerate}
\end{sol}